% Generated draft: requires mathematical and source review.
\documentclass[11pt]{article}
\usepackage[margin=25mm]{geometry}
\usepackage{fontspec}
\setmainfont{Latin Modern Roman}
\usepackage{amsmath,amssymb,mathtools,unicode-math}
\setmathfont{Latin Modern Math}
\usepackage{enumitem,needspace,xurl,hyperref}
\setlength{\parindent}{0pt}
\setlength{\parskip}{5pt}
\setlength{\emergencystretch}{3em}
\newcommand{\N}{\mathbb{N}}
\newcommand{\Z}{\mathbb{Z}}
\newcommand{\Q}{\mathbb{Q}}
\newcommand{\R}{\mathbb{R}}
\newcommand{\C}{\mathbb{C}}
\newcommand{\F}{\mathbb{F}}
\newcommand{\A}{\mathbb{A}}
\newcommand{\PP}{\mathbb{P}}
\newcommand{\E}{\mathbb{E}}
\newcommand{\eps}{\varepsilon}
\newcommand{\dd}{\,\mathrm{d}}
\begin{document}
\section*{UnsolvedMath ID: 30000192}
\textbf{Source problem code:} catalogue entry
\subsection{Definitions and mathematical statement}
\textbf{Research question.} Does a MacWilliams duality theorem hold for general classes of convolutional codes? The source asks for an extension of a classical duality and explicitly expresses belief that it holds for much more general classes, leaving that extension for future research. It does not define the target class, code dual, or enumerator in the surviving statement. The attached assessment distinguishes ordinary weight distributions, for which it reports failure of a general convolutional-code duality, from a weight-adjacency-matrix duality reported for degree-one codes. These are different target assertions.\par A conventional algebraic model is a rank-\(r\) convolutional code \(\mathcal C\subset\F_q[z]^n\), represented by a polynomial encoder and, in the intended class, a finite-state realization. Choices such as basic or noncatastrophic encoders, polynomial versus rational sequence spaces, and the bilinear form or time-reversal convention defining \(\mathcal C^\perp\) must be specified. For a chosen minimal state realization, a weight adjacency matrix has entries \(\Lambda_{s,t}(W)=\sum_e W^{\operatorname{wt}(v_e)}\), summing over allowed one-step transitions from state \(s\) to state \(t\), where \(v_e\in\F_q^n\) is the output and \(\operatorname{wt}\) is Hamming weight. Multiplicities and state-identification conventions belong to that definition. A MacWilliams-type statement in this setting should determine the corresponding invariant of \(\mathcal C^\perp\) by an explicitly specified character transform and weight-variable transformation of the invariant of \(\mathcal C\), with stated normalizations and equivalence under permitted state-coordinate changes. No such complete formula is supplied here.\par The surviving source display gives \(H_{x,y}=\chi(xy)\), \(x,y\in\F_q\), with \(\chi:(\F_q,+)\to\C^\times\) a nontrivial additive character, and contains a detached fragment \(1+(q-1)W\). The full matrix identity, any factors of \(q\), substitutions, and state-space extension are missing; the fragment is not promoted to a recovered theorem. The code degree in the reported special case refers to the minimal state-space degree, not automatically to block length or rate.\par The supplied literature triage checked 2026-08-21 labels the problem partially solved, reporting ordinary-weight-distribution failure, a degree-one weight-adjacency theorem, and an unresolved broader extension. The precise special-case hypotheses and current scope of any later theorem are not reproduced. This dated assessment requires independent verification and is not adopted as a newly checked status.
\subsection{Short English statement}
Seek a MacWilliams-type transform relating suitable convolutional-code invariants to those of their duals. The intended generalization concerns richer state-based data rather than an ordinary weight-distribution identity already reported to fail. The exact code class, duality convention, and transform formula are missing.
\subsection{Sources}
\begin{itemize}\item Local UnsolvedMath registry, permanent ID 30000192, code catalogue entry; \emph{MacWilliams Duality for Convolutional Codes}, Oberwolfach Reports Open Problems.\item \emph{Regelungstheorie (2005)}, Oberwolfach Reports, DOI 10.4171/owr/2005/11, pp. 565--645. \url{https://doi.org/10.4171/owr/2005/11}. Proposed year 2005; no proposer named.\item Supplied literature label \emph{original problem and cited results}: \url{https://doi.org/10.4171/owr/2005/11}.\item Supplied literature label \emph{report text with the degree-one theorem}: \url{https://ems.press/content/serial-article-files/45984?nt=1}. No separate theorem citation or complete displayed identity is supplied.\item Literature triage checked 2026-08-21, status \texttt{partially\_solved}; record updated 2026-09-16. The \texttt{corrected\_verified} field records removal of expository spillover and a corrupted matrix display; dated claims were not independently checked here.\end{itemize}
\end{document}
